% !TeX program = pdflatex
\documentclass[UTF8,a4paper,11pt]{ctexart}

\usepackage[a4paper,margin=2.2cm]{geometry}
\usepackage{booktabs}
\usepackage{hyperref}
\usepackage{enumitem}
\usepackage{listings}
\usepackage{amsmath,amssymb}
\usepackage{graphicx}
\usepackage{caption}
\usepackage{longtable}
\usepackage{array}

\hypersetup{
  colorlinks=true,
  linkcolor=blue,
  urlcolor=blue,
  citecolor=blue
}

\lstset{
  basicstyle=\ttfamily\small,
  breaklines=true,
  frame=single,
  columns=fullflexible
}

\title{E--PIC 代码结构重组与多物理模块化方案（mesh.dict/geometry.dict 统一驱动）}
\author{方案整理：ChatGPT（基于 E-PIC-Task3 代码快照）}
\date{2026-01-21}

\begin{document}
\maketitle

\begin{abstract}
本文给出一份面向 E--PIC 代码的“结构重组 + 可拓展多物理模块化 + 便捷算例管理”的工程方案。
核心目标是：\textbf{(1) 网格与几何统一由 \texttt{mesh.dict} 与 \texttt{geometry.dict} 驱动；(2) 物理功能以可插拔 \texttt{Module} 形式组织，\texttt{main()} 只保留最小启动逻辑；(3) 算例以子目录形式管理，运行时切换 case 而非覆盖根目录 \texttt{system}。}
方案在尽量少破坏现有代码的前提下，提出分阶段迁移路径、目录结构、字典组织、运行流程与扩展建议。
\end{abstract}

\tableofcontents
\vspace{0.8em}

%======================================================================
\section{背景与需求梳理}

用户需求（归纳）：
\begin{enumerate}[label=\arabic*)]
  \item 以统一的几何/网格定义（\texttt{geometry.dict}, \texttt{mesh.dict}）作为\emph{所有}物理模块的基础输入；
  \item 将 Electricity、Hydrodynamics、ExplicitPIC 等作为可独立拓展的\textbf{功能模块}（Module），未来可继续增加新的物理/数值模块；
  \item \texttt{main()} 文件“写薄”：把当前主程序里与电喷雾相关的具体流程抽到\textbf{电喷雾模块}，像 ExplicitPICModule 一样通过 \texttt{RunController} 入口切换；
  \item 删除 \texttt{Physics/} 目录：其中内容合并到 \texttt{Modules/}；\texttt{ImmersedGeometry} 移入 \texttt{Geometry/} 作为基础能力；
  \item 算例管理：不再每次覆盖根目录 \texttt{system/}；希望每个算例一个子目录，运行时切换（类似 OpenFOAM 的 case 组织）。
\end{enumerate}

本文方案额外强调：\textbf{保持“核心库 + 模块插件 + case”三层解耦}，使得后续增加模块、替换求解器或做多物理耦合时，最少改动、不互相污染。

%======================================================================
\section{现状快照（基于 E-PIC-Task3 代码结构）}

当前目录结构（简化）：
\begin{lstlisting}
include/
  Base/ Discrete/ Field/ Geometry/ IO/ Lagrangian/ Mesh/
  Modules/ExplicitPIC/...
  Physics/ Electricity.h Hydrodynamics.h ImmersedGeometry.h ...
  Runtime/ IModule.h RunController.h
src/
  ...
  Main.cpp
  Physics/ Electricity.cpp Hydrodynamics.cpp
  Runtime/ RunController.cpp
system/
  mesh.dict geometry.dict run.dict pic.dict ...
\end{lstlisting}

当前已具备的“模块化雏形”：
\begin{itemize}
  \item \texttt{Runtime/IModule}：轻量模块接口；
  \item \texttt{Runtime/RunController}：读取 \texttt{./system/run.dict}，目前只识别 \texttt{ExplicitPIC}，否则进入 legacy 的巨大 \texttt{Main.cpp}；
  \item \texttt{DictionaryRun/DictionaryPIC} 已支持按路径读取；但 \texttt{DictionaryMesh/DictionaryGeometry} 仍固定读取 \texttt{./system/*.dict}。
\end{itemize}

主要痛点：
\begin{enumerate}[label=\arabic*)]
  \item \textbf{主程序过厚}：电喷雾/电流体等具体流程与底层库混在 \texttt{Main.cpp}；
  \item \textbf{目录语义不清}：\texttt{Physics/} 同时承载“可复用的物理场求解器类”和“某些模块化入口”的潜在角色；
  \item \textbf{字典路径硬编码}：Mesh/Geometry 的字典只能从根目录 \texttt{system/} 读，不支持按 case 切换；
  \item \textbf{运行组织能力弱}：无法声明“多模块耦合/顺序执行/共享网格几何/共享输出”等。
\end{enumerate}

%======================================================================
\section{目标架构：Core + Geometry/Mesh + Modules + Case}

\subsection{总体分层}
\begin{itemize}
  \item \textbf{Core（基础库）}：Base/IO/Field/Discrete/Lagrangian/Mesh 等，提供通用数据结构与数值工具。
  \item \textbf{Geometry/Mesh 基础能力}：几何描述、水平集、EB/浸入边界几何（\texttt{ImmersedGeometry}）等，\emph{不与具体物理强耦合}。
  \item \textbf{Modules（功能模块）}：每个模块是一个可独立运行的入口（实现 \texttt{IModule}），内部可复用 Core 的类与算法。
  \item \textbf{Cases（算例）}：每个算例一个目录，包含 \texttt{system/} 子目录，存放 \texttt{mesh.dict, geometry.dict, run.dict} 及各模块字典。
\end{itemize}

\subsection{关键原则}
\begin{enumerate}[label=\arabic*)]
  \item \textbf{网格/几何单一来源}：\texttt{mesh.dict} 与 \texttt{geometry.dict} 由 Runtime 在启动时统一读取并构建 Mesh/Geometry；模块不得各自重复读一套“私有网格几何”。
  \item \textbf{模块配置隔离}：每个模块只读取自己的 dict（例如 \texttt{electricity.dict}、\texttt{hydrodynamics.dict}、\texttt{pic.dict}、\texttt{electrospray.dict}），避免“全局巨 dict”。
  \item \textbf{main 只做启动}：解析命令行/确定 case 路径/初始化 RunController/选择模块并运行。
  \item \textbf{模块可组合}：电喷雾模块应作为 driver（驱动器模块）去组合 Electricity/Hydrodynamics/ExplicitPIC 等能力。
\end{enumerate}

%======================================================================
\section{建议的目录结构重组}

\subsection{建议的目标目录树}
\begin{lstlisting}
E-PIC/
  CMakeLists.txt
  apps/
    epic_main.cpp                 # 极薄 main
  include/
    Core/...                      # Base/IO/Field/Discrete/Lagrangian/Mesh
    Geometry/
      LevelSet.* Line.* Phase.*
      ImmersedGeometry.*          # 从 Physics 迁入
    Runtime/
      IModule.h
      RunController.*
      Case.h Case.cpp             # case 路径/字典定位
      ModuleRegistry.*            # 模块注册/工厂
      SimContext.*                # 共享上下文（Mesh/Geometry/输出等）
    Modules/
      ExplicitPIC/
      Electricity/
      Hydrodynamics/
      Electrospray/               # 新：把 Main.cpp 的电喷雾流程迁入
  src/
    ... 与 include 对应实现 ...
  cases/
    spray2D/
      system/
        run.dict mesh.dict geometry.dict
        electricity.dict hydrodynamics.dict pic.dict electrospray.dict
      constant/ (可选)
      0/ (可选)
\end{lstlisting}

\subsection{Physics/ 的处理}
\begin{itemize}
  \item 将 \texttt{include/Physics/Electricity.*} 迁入 \texttt{include/Modules/Electricity/} 或 \texttt{include/Models/Electricity/}。
  \item 将 \texttt{Hydrodynamics.*} 同理迁入。
  \item \texttt{ImmersedGeometry.*} 迁入 \texttt{include/Geometry/}（建议保持类名不变以减少改动）。
  \item \texttt{Physics/} 目录删除。
\end{itemize}

关于“Electricity/Hydrodynamics 是模块还是可复用 solver 类”的建议：
\begin{itemize}
  \item 保留它们为\textbf{可复用的 solver/model 类}（例如 \texttt{ElectricitySolver}），同时在 \texttt{Modules/Electricity} 下提供一个轻量入口模块 \texttt{ElectricityModule}。
  \item 这样既满足“模块化运行”，又避免把求解器类写成只能单独跑的脚本式程序。
\end{itemize}

%======================================================================
\section{Case 组织与字典定位（解决覆盖 system 的痛点）}

\subsection{Case 目录约定}
建议采用 OpenFOAM 风格最小集合：
\begin{itemize}
  \item \texttt{<case>/system/}：运行与离散相关字典（网格/几何/模块参数）；
  \item \texttt{<case>/output/}：统一输出；
  \item 未来可扩展：\texttt{constant/}（物性/常数）、\texttt{0/}（初场）。
\end{itemize}

\subsection{新增 Runtime::Case（字典路径服务）}
新增一个轻量对象 \texttt{Case}：
\begin{lstlisting}
struct Case {
  std::filesystem::path root;          // case 根目录
  std::filesystem::path systemDir();   // root/system
  std::filesystem::path dict(std::string name);
};
// 用法：case.dict("mesh.dict") -> <case>/system/mesh.dict
\end{lstlisting}

\subsection{命令行接口（推荐）}
建议支持：
\begin{lstlisting}
# 默认：--case ./cases/default
./epic --case ./cases/spray2D

# 允许覆盖 run.dict 的 module
./epic --case ./cases/spray2D --module Electrospray
\end{lstlisting}

\subsection{对现有 Dictionary* 的最小侵入改造}
\begin{itemize}
  \item \texttt{DictionaryMesh::readDict()} 改为 \texttt{readDict(path)} 或构造时传入 \texttt{Case}；
  \item \texttt{DictionaryGeometry} 同理；
  \item \texttt{DictionaryRun/DictionaryPIC} 已支持传路径，可直接复用；
  \item 保留默认路径以兼容旧用法：若未传 case，则仍读 \texttt{./system/*.dict}。
\end{itemize}

%======================================================================
\section{Runtime 与 Modules：从“单模块入口”走向“多物理可组合”}

\subsection{模块注册与工厂（避免 if-else 链）}
当前 \texttt{RunController} 用 if-else 选择模块，模块多了会膨胀。
建议引入 \texttt{ModuleRegistry}：
\begin{lstlisting}
using ModuleFactory = std::function<std::unique_ptr<IModule>(const RunConfig&, SimContext&)>;
class ModuleRegistry {
public:
  static ModuleRegistry& instance();
  void add(std::string name, ModuleFactory f);
  std::unique_ptr<IModule> create(std::string name, const RunConfig&, SimContext&);
};
// 每个模块 cpp 内 static 注册：REGISTER_MODULE("ExplicitPIC", ExplicitPICModule);
\end{lstlisting}

\subsection{SimContext：共享网格/几何/输出/随机数/并行环境}
为满足“网格几何统一构建一次，模块共享”，建议在 Runtime 层建立 \texttt{SimContext}：
\begin{lstlisting}
struct SimContext {
  Case theCase;
  std::unique_ptr<QuadtreeMesh> mesh;
  std::unique_ptr<ImmersedGeometry<QuadtreeMesh>> iGeom;
  OutputManager out;      // 统一输出（目录、命名、步号）
  Logger log;             // 统一日志
  // 可选：random、MPI/OpenMP 信息、profiling、global parameters...
};
\end{lstlisting}

构建流程：
\begin{enumerate}[label=\arabic*)]
  \item RunController 解析 case；
  \item 读取 \texttt{mesh.dict} 构建 mesh；读取 \texttt{geometry.dict} 构建 level set / lines / phases；
  \item 构建 \texttt{ImmersedGeometry(mesh, geometry)}；
  \item 再创建并运行选定 Module（Module 仅从 context 拿 mesh/iGeom）。
\end{enumerate}

\subsection{模块生命周期建议（为未来多物理耦合预留）}
当前 \texttt{IModule::run()} 够用，但为了后续“耦合/分裂/多步驱动”，建议逐步扩展为：
\begin{lstlisting}
class IModule {
public:
  virtual void initialize() {}
  virtual void step(int n, double t, double dt) = 0;
  virtual void finalize() {}
};
\end{lstlisting}

说明：不必一次改完；可以先在 \texttt{ElectrosprayModule} 内部实现自己的 step loop，并保持其他模块仍是 run()。

%======================================================================
\section{电喷雾模块化落地：从 Main.cpp 迁移}

\subsection{目标：把“电喷雾程序”变成 ElectrosprayModule}
建议新建：\texttt{Modules/Electrospray/ElectrosprayModule.*}
\begin{itemize}
  \item \textbf{输入字典}：\texttt{electrospray.dict}（耦合策略、宏步长、粒子注入、喷射模型参数等）；
  \item \textbf{依赖服务}：从 \texttt{SimContext} 获取 mesh/iGeom；内部组合 Electricity/Hydrodynamics/ExplicitPIC 的求解器对象。
  \item \textbf{输出}：统一走 OutputManager，避免各处硬编码 \texttt{./output}。
\end{itemize}

\subsection{Main.cpp 中应迁移的典型内容}
从当前 \texttt{Main.cpp} 可见：时间步计算、输出时间序列、体积分数/界面体积统计、以及电喷雾耦合流程。
迁移建议：
\begin{enumerate}[label=\arabic*)]
  \item 把电喷雾专属函数（如 \texttt{compute\_dt}, \texttt{outputTemporalResults}）移入 ElectrosprayModule 内部（或其 \texttt{detail} 命名空间）；
  \item 把所有硬编码输出路径改成 \texttt{context.out.path("time\_results.dat")} 之类；
  \item 把模型常数（\texttt{character\_rho} 等）从代码常量迁移到 \texttt{electrospray.dict} 或 \texttt{constant/}。
\end{enumerate}

\subsection{ElectrosprayModule 与子求解器的关系}
推荐把 Electricity/Hydrodynamics/PIC 当作\textbf{可复用求解器类}：
\begin{itemize}
  \item \texttt{ElectricitySolver}：提供 assemble/solve/gather 等；
  \item \texttt{HydrodynamicsSolver}：提供压力速度耦合、界面重构等；
  \item \texttt{ExplicitPICDriver}：显式 PIC 的推进与沉积。
\end{itemize}

ElectrosprayModule 作为 driver：
\begin{lstlisting}
initialize():
  build fields, read dicts
for n=0..N-1:
  (1) hydrodynamics predictor/corrector
  (2) electricity solve (with updated interface)
  (3) PIC advance / injection / collisions
  (4) diagnostics + output
finalize()
\end{lstlisting}

这样既满足“电喷雾是一个入口模块”，又保留底层求解器可被其他 driver 复用（例如“等离子体放电模块”“介质界面模块”等）。

%======================================================================
\section{Run.dict 的建议扩展（支持 case 切换与多模块）}

保持现有风格，建议在 \texttt{run.dict} 支持：
\begin{lstlisting}
Run
{
  module   Electrospray
  # 可选：case 由命令行指定；也可在此声明
  # case    ./cases/spray2D
  logLevel info
}

Modules
{
  Electrospray { dict electrospray.dict }
  Electricity  { dict electricity.dict  }
  Hydro        { dict hydrodynamics.dict }
  PIC          { dict pic.dict }
}
\end{lstlisting}

说明：\texttt{Modules} 段仅用于声明“每个模块的配置文件名”，避免把所有参数都写进 run.dict。

%======================================================================
\section{迁移实施路线（推荐分 4 个阶段）}

\subsection{Phase 0：不破坏旧功能的前提下引入 case}
\begin{enumerate}[label=\arabic*)]
  \item 新增 \texttt{Runtime::Case}；\texttt{RunController} 增加命令行解析（\texttt{--case}, \texttt{--module}）。
  \item 改造 \texttt{DictionaryMesh/DictionaryGeometry} 支持按路径读取。
  \item 若未提供 case：默认仍用 \texttt{./system}，确保旧用法不崩。
\end{enumerate}

\subsection{Phase 1：ImmersedGeometry 迁移到 Geometry}
\begin{enumerate}[label=\arabic*)]
  \item 移动头文件与实现到 \texttt{Geometry/}；
  \item 修正 include 路径与 CMake；
  \item 对外接口保持不变（类名/函数签名尽量不动）。
\end{enumerate}

\subsection{Phase 2：Physics 拆解并形成 Electricity/Hydro 模块}
\begin{enumerate}[label=\arabic*)]
  \item 把 \texttt{Electricity.* / Hydrodynamics.*} 迁入 \texttt{Modules/}（或 \texttt{Models/} + \texttt{Modules/} 双层）；
  \item 增加 \texttt{ElectricityModule} 与 \texttt{HydrodynamicsModule} 的最小入口（可先只做 dry-run/单步求解/输出验证）；
  \item RunController 通过 ModuleRegistry 支持多模块选择。
\end{enumerate}

\subsection{Phase 3：电喷雾模块落地，main 变薄}
\begin{enumerate}[label=\arabic*)]
  \item 新建 \texttt{ElectrosprayModule}，把 \texttt{Main.cpp} 里电喷雾相关流程迁入；
  \item \texttt{apps/epic\_main.cpp} 只保留：初始化、选择 case、创建 context、run 模块；
  \item legacy Main.cpp 可暂时保留为 \texttt{LegacyElectrospray} 模块或直接删除（建议先保留一段时间便于回归对比）。
\end{enumerate}

%======================================================================
\section{除了你提出的点以外，额外的可拓展/易用性建议}

\subsection{建议 1：统一输出与目录管理（OutputManager）}
现有代码中输出路径多处硬编码 \texttt{./output}。
建议新增 \texttt{OutputManager}：
\begin{itemize}
  \item 所有输出都写到 \texttt{<case>/output/}；
  \item 支持按步号自动建子目录（\texttt{output/000100/}），避免文件互相覆盖；
  \item 为后续并行（MPI rank）输出预留前缀。
\end{itemize}

\subsection{建议 2：字典 include 与默认配置（减少重复）}
很多 case 之间只有少量参数差异。
建议字典支持：
\begin{itemize}
  \item \texttt{include "../defaults/mesh.dict"} 之类；
  \item 变量替换（例如 \texttt{dt = \$\{dt\_macro\}}）。
\end{itemize}

若暂时不想增强解析器，也可通过“默认模板复制 + override 小文件”的方式实现。

\subsection{建议 3：模块间依赖声明与运行管线}
多物理 driver 之外，有时希望顺序运行多个模块（例如先生成几何，再求电，再跑 PIC）。
建议 run.dict 支持：
\begin{lstlisting}
Run { pipeline ( GeometryCheck Electricity PIC ) }
\end{lstlisting}

并由 RunController 循环创建并运行。

\subsection{建议 4：回归测试与最小算例集}
重构期最怕“改完能编译但结果悄悄变了”。建议建立：
\begin{itemize}
  \item \textbf{unit tests}：字典解析、网格生成、EB 几何度量等；
  \item \textbf{mini cases}：Poisson/Laplace、静态界面几何、PIC charge deposition 守恒等；
  \item 在 CMake 里加入 \texttt{ctest}，每次改动自动跑。
\end{itemize}

\subsection{建议 5：明确命名与 API 边界}
建议约定：
\begin{itemize}
  \item \texttt{*Solver}：可复用求解器类（不含主循环）；
  \item \texttt{*Module}：可独立运行入口（含主循环或 driver）；
  \item \texttt{Dictionary*}：只做配置解析，不承载业务逻辑。
\end{itemize}

%======================================================================
\section{风险与对策}

\begin{longtable}{p{0.26\textwidth}p{0.34\textwidth}p{0.34\textwidth}}
\toprule
风险 & 典型表现 & 对策\\
\midrule
重构期间功能回归 & 编译过但结果变、输出丢失 & Phase 0 起就建立 mini case + 输出对比；保留 legacy 模式一段时间\\
字典体系膨胀 & 参数分散、难以定位 & 约定每模块一个 dict；run.dict 只做选择与引用；提供 defaults/include\\
模块间共享数据混乱 & 字段重复、生命周期不清 & SimContext 统一托管 mesh/geometry/output；模块只拿引用\\
未来并行/性能改动困难 & 日志/输出/随机数各自为政 & 统一 OutputManager/Logger/Random；预留 MPI/OpenMP 信息\\
\bottomrule
\end{longtable}

%======================================================================
\section{附录：建议的最小字典模板}

\subsection*{A. run.dict}
\begin{lstlisting}
Run
{
  module   Electrospray
  logLevel info
}

Modules
{
  Electrospray { dict electrospray.dict }
  Electricity  { dict electricity.dict }
  Hydro        { dict hydrodynamics.dict }
  PIC          { dict pic.dict }
}
\end{lstlisting}

\subsection*{B. mesh.dict / geometry.dict}
保持你现有格式不变，只把文件位置挪到 \texttt{<case>/system/}。

\subsection*{C. electrospray.dict（示意）}
\begin{lstlisting}
Electrospray
{
  dtMacro   1e-8
  steps     10000
  coupling  semiImplicit
  outputEvery 50
  # model constants / scaling...
}
\end{lstlisting}

\end{document}
